\section{Preliminaries}\label{prelim} 

Throughout we set $[n]\coloneqq \{1, \ldots , n\}$ and we use $\mathbf{I}_{k, n}$ to denote the collection of subsets of $[n]$ of size $k$. The symmetric group on $k$ elements is denoted by $S_k$. We also fix a field $\mathbb{K}$. 

\smallskip

The Grassmannian $\Gr(k, n)$ is the space of all $k$ dimensional linear subspaces of $\mathbb{K}^n$. A point in $\Gr(k, n)$ can be represented by a $k\times n$ matrix with entries in $\mathbb{K}$. 
Let $X=(x_{ij})$ be a $k\times n$ matrix of indeterminates. For a subset $I = \{i_1,\ldots,i_k\} \in \mathbf{I}_{k, n}$, let $X_I$ denote the $k\times k$ submatrix with the column indices $i_1,\ldots,i_k$. 
The \emph{Pl{\"u}cker forms (or Pl{\"u}cker coordinates)} are $P_I = \det (X_I)$ for $I \in\mathbf{I}_{k, n}$. These forms determine the \emph{Pl{\"u}cker embedding} from
$ \Gr (k,n)$ into $\mathbb{P}^{\binom{n}{k}-1}$. 

In the following, we consider the polynomial ring $\mathbb{K}[x_{ij}]$ on the variables $x_{ij}$ with $1\leq i\leq k$ and $1\leq j\leq n$ and the polynomial ring $\mathbb{K}[P_I]$ on the Pl{\"u}cker variables with $|I|=k$.

\begin{Definition}\label{def:ideal}The Pl{\"u}cker ideal $\I_{k,n}$ is defined as the kernel of the map 
\begin{equation}\label{eqn:pluckermap}
\begin{aligned} 
\psi \colon \mathbb{K}[P_I] &\rightarrow \mathbb{K}[x_{ij}] 
\\ 
 P_{I} &\mapsto \det (X_I).
\end{aligned}
\end{equation}
The \emph{Pl{\"u}cker algebra} is the finitely generated algebra $\mathbb{K}[P_I ]/{\mathcal{I}_{k,n}}$ 
denoted by $\mathcal{A}_{k,n}$.
\end{Definition}




\begin{Definition}\label{def:matching}
A $k\times n$ matching field is a map $\Lambda: \mathbf{I}_{k,n} \to S_k$.
\end{Definition}

Given a $k\times n$ matching field $\Lambda$ and a subset $I = \{i_1, \ldots , i_k\} \in \mathbf{I}_{k, n}$ 
we consider the set to be ordered by $i_1 < \cdots < i_{k}$. We think of the permutation {$\sigma=\Lambda(I)$} as inducing a new ordering on the elements of $I$ 
where the position of $i_s$ is determined by the value of $\sigma(s)$.
It is convenient to represent the variable $P_I$ as a $k \times 1$ tableau 
 where $(\sigma(r), 1)$ contains $i_{r}$.

\begin{Definition}
Let $\Lambda$ be a size $k \times n$ matching field. 
A $\Lambda$-tableau is a tableau of size $k \times d$ for any $d\geq 1$ with entries in $[n]$,
so that the entries in each column are pairwise distinct and filled according to the order determined by $\Lambda$. 
\end{Definition}


\begin{exam}\label{ex:diagonal1}
The \emph{diagonal matching field} assigns to
each subset $I \in \mathbf{I}_{k,n}$ the identity permutation~\cite[Example~1.3]{sturmfels1993maximal}.
Therefore, a $\Lambda$-tableau is a rectangular tableau of size $k \times d$ filled with entries in $[n]$ such that the columns are strictly increasing. 
\end{exam}

\begin{exam}\label{ex:pointed1}
A $k \times n$ matching field is called \emph{pointed} if there exists $i_1, \ldots, i_k \in [n]$ such that if $i_s \in I$ for some 
$1 \leq s \leq k$ then $\Lambda(I)(s) = s$~\cite[Example~1.4]{sturmfels1993maximal}. In other words, if $i_s \in I$ for some $1\leq s \leq k$ then the column corresponding to $P_I$ contains $i_s$ in row $s$. Below are the tableaux representing $P_I$ for a pointed matching field of size $3 \times 5$ which is otherwise filled diagonally:
\[
\begin{array}{c}1 \\2 \\ 3 \end{array} , \quad
\begin{array}{c} 1 \\2 \\4 \end{array},\quad
\begin{array}{c}1 \\2 \\ 5 \end{array} , \quad
\begin{array}{c}1 \\ 4 \\ 3 \end{array} ,\quad
\begin{array}{c}1 \\5 \\ 3\end{array} ,\quad
\begin{array}{c}4 \\ 2\\ 3 \end{array} ,\quad
\begin{array}{c}5 \\2 \\ 3\end{array} ,\quad
\begin{array}{c} 4 \\ 2 \\ 5 \end{array}. 
\]
Generalising this, we say that a size $k \times n$ matching field $\Lambda$ is pointed on $S \subset [n]$ if for all $i \in S$ there exists a $j_i$ such that $i$ always appears in row $j_i$ of any $\Lambda$ matching field tableau. 
Here $S$ need not have size equal to $k$. 
\end{exam}


A monomial $\Pi_{I \in A} P_I $ can be represented by a $\Lambda$-tableau of size $k \times |A|$ given by the concatenation of the columns with content $I$ filled according to the matching field. 
To each monomial $\Pi_{I \in A} P_I $ we associate a sign $\epsilon_{A} = \pm 1$ determined by the signature of the permutations $\Lambda(I)$ for all $I \in A$. More precisely, 
\[
\epsilon_{A} = \Pi_{I \in A} \sgn (\Lambda(I)).
\]


\begin{Definition}
Given a matching field $\Lambda$, the \emph{matching field ideal} $\mathcal{I}_{\Lambda} \subset \mathbb{K}[x_{ij}]$ is generated by the 
 binomial relations 
\begin{equation}\label{eq:binomials}
\epsilon_A \Pi_{I \in A} P_I - \epsilon_B \Pi_{J \in B} P_J
\end{equation}
if and only if the contents of the corresponding $\Lambda$-tableau of size $k \times |A|$ are row-wise equal.
\end{Definition}

To $I \in \mathbf{I}_{k, n}$ {with $\sigma=\Lambda(I)$} we associate the monomial 
\[ 
\mathbf{x}_{\Lambda(I)}\coloneqq x_{\sigma(1) i_{1}}x_{\sigma(2)i_2}\cdots x_{{\sigma(k)i_k}}. 
 \]
A $k\times n$ matching field $\Lambda$, 
gives a map of polynomial rings 
\begin{equation}\label{eqn:monomialmap}
\begin{aligned}
\phi_{\Lambda} \colon\ \mathbb{K}[P_I] &\rightarrow \mathbb{K}[x_{ij}] 
\\ 
 P_{I} &\mapsto \sgn(\Lambda(I)) \mathbf{x}_{\Lambda(I)}.
\end{aligned}
\end{equation}

\smallskip



\begin{Proposition}
Given a matching field $\Lambda$, the matching field ideal $\mathcal{I}_{\Lambda}$ is the kernel of the monomial map
$\phi_{\Lambda}$ from Equation~\eqref{eqn:monomialmap}. 
\end{Proposition}


